\documentclass[10pt,a4paper]{amsart}
\usepackage[margin=19mm]{geometry}
\usepackage{amsmath,amssymb,mathtools}
\usepackage[scr=rsfs,cal=euler]{mathalfa}
\usepackage{xfrac}
\usepackage[unicode,hidelinks,bookmarks=false]{hyperref}
\newcommand{\ie}{i.e.\ }
\newcommand{\cf}{cf.\ }
\newcommand{\resp}{respectively\ }
\newcommand{\Subsection}{\subsection}
\providecommand{\nobreakdash}{\nobreak\hbox{-}}
\setlength{\emergencystretch}{2em}
\multlinegap0pt
\usepackage{bm}
\usepackage{bbm}
\usepackage{dsfont}
\usepackage{xspace}
\usepackage{cancel}
\usepackage{mathtools}
\usepackage{amsthm}
\makeatletter
\@ifundefined{theorem}{\newtheorem{theorem}{Theorem}}{}
\@ifundefined{claim}{\newtheorem{claim}[theorem]{Claim}}{}
\makeatother
\def\mD{\mathbb{D}}
\newcommand{\Smax}[1][r]{\Sigma_{\max}(#1)}
\newcommand{\Smin}[1][r]{\Sigma_{\min}(#1)}
\newcommand{\amax}[1][r]{x_{\max}(#1)}
\newcommand{\amin}[1][r]{x_{\min}(#1)}
\newcommand{\ep}{\varepsilon}
\newcommand{\pppp}[1]{\partial^{++}(#1)}
\title[Any fully graphic region of degree sequences can be sampled ]{Any fully graphic region of degree sequences can be sampled rapidly}
\author{P\'eter L. Erd\H{o}s, G\'abor Lippner, Na'ama Nevo, Lajos Soukup}
\date{Source: arXiv:2511.13564v1 (CC BY 4.0)}
\begin{document}
\maketitle

\noindent\emph{Source and licence.} Any fully graphic region of degree sequences can be sampled rapidly. Version arXiv:2511.13564v1 (CC BY 4.0), distributed under the Creative Commons Attribution 4.0 International licence, as recorded in the arXiv licence metadata for this exact version and shown on the abstract page https://arxiv.org/abs/2511.13564v1. Full rights notice: licenses/arxiv-2511.13564-CC-BY-4.0.txt.\par\smallskip

\noindent\textit{Editorial scope.} Counted unit: the Claim whose source label is \texttt{eq:newxbound} in the article (source0.tex lines 1329--1336, source label \texttt{eq:newxbound}), delivered together with the article's own complete original proof of it (source0.tex lines 1353--1412, one \texttt{proof} environment of 60 lines). No mathematics is written, repaired or re-derived: statement and proof are the article's own text line for line. The proof is detached from the statement in the source: the article states the result at lines 1329--1336 and prints its proof at lines 1353--1412, and the proof environment's own header names the statement, so the association is the article's own. Required context retained uncounted: thm:sigmabound (lines 266--283); eq:sigma\_abs\_bound (lines 274--277). Every definition, display and label of those lines is retained unchanged. Omitted from this package and not redistributed: 1--265, 278--1328, 1337--1352, 1413--1460. Nothing inside a retained excerpt is omitted or altered: every retained body is delivered exactly as the article's lines are written, so no omission receipt of any kind accompanies this package. Citations: the retained excerpts carry no citation command at all, so no bibliography entry of the article is transcribed and no reference is redistributed. Reference adaptation: none was needed and none was made; every reference inside the delivered unit resolves inside it, so no unresolved reference or citation is produced. Printed slips kept literal: no printed word, symbol, hypothesis or number of the counted statement or of its proof was altered. Layout and class: the article's own class is \texttt{elsarticle} with packages including geometry, times, amsthm, amsmath, amssymb, mathtools, relsize, fontenc, inputenc, enumerate, float, verbatim, hyperref, csquotes; the wrapper is the lane's pinned \texttt{amsart} editorial wrapper at 10pt on A4, which restates the theorem-like environments the retained text needs and adds the article's own macro definitions that the retained text needs (\texttt{\textbackslash Smax}, \texttt{\textbackslash Smin}, \texttt{\textbackslash amax}, \texttt{\textbackslash amin}, \texttt{\textbackslash ep}, \texttt{\textbackslash mD}, \texttt{\textbackslash pppp}), with extra wrapper packages (bm, bbm, dsfont, xspace, cancel, mathtools). The delivered numbering is the unit's own. Assets: no retained span contains a figure environment, a graphics-inclusion command, a PDF-page inclusion, a table or a TikZ picture; no image file is copied into this package, the package declares no asset dependency and no image file is redistributed with it. Licence: version arXiv:2511.13564v1 of this article is distributed under the Creative Commons Attribution 4.0 International licence, as recorded in the arXiv licence metadata for this exact version and shown on the abstract page https://arxiv.org/abs/2511.13564v1. A full rights notice is delivered at licenses/arxiv-2511.13564-CC-BY-4.0.txt. Characters: every retained excerpt is pure ASCII, so the delivered engine, XeLaTeX with its default Latin Modern text font, reports no missing character.
\begin{theorem}\label{thm:sigmabound}
    Suppose $\beta > 0$,  $r\in \mathbb N$ and
    $\varepsilon = \tfrac{3(r+3)}{\beta(c_1 - c_2)}$.
    If
    \begin{equation}\label{eq:beta}
        (1-\beta)(c_1 - c_2 + 1)^2 \geq 4c_2(n - 1 - c_1).
    \end{equation}
    and
    \begin{equation}\label{eq:sigma_abs_bound}
    \left| \frac{\Sigma - n c_2}{(c_1 - c_2)/2} - (1 + c_1 + c_2) \right|
        \leq (1 - \varepsilon)\sqrt{Q},
    \end{equation}
    then
     the simple region $\mD(n,\Sigma,c_1,c_2)$ contains a sequence $D$ such that
    \begin{equation}\label{eq:notstable}
       \pppp D \ge  2^{r/2}.
        \end{equation}
        \end{theorem}

\begin{claim}
    $I^x \cap I^{x+1} \neq \emptyset$ if and only if $I^{x+1}_0 \leq I^{x}_1$, that is when
    \[ (x+1)(x+2r)+2c_2(n-x-1-r) \leq x(x+2r-1) + 2x(c_1-x-r-+1),\]
    or, equivalently,
    \begin{equation}\label{eq:newxbound}
        x^2-x(c_1+c_2-r)+r + c_2(n-1-r) \leq 0
    \end{equation}
\end{claim}

\begin{proof}[Proof of Theorem~\ref{thm:sigmabound}] All we need to establish that the interval defined by \eqref{eq:sigma_abs_bound} is contained in $[\Smin,\Smax]$. Or rather, we need to determine a value of $\ep$ that allows us to establish this, and then we need to prove that this $\ep \leq \frac{3(r+3)}{\beta(c_1-c_2)}$.


Using that $\amin$ and $\amax$ satisfy \eqref{eq:newxbound} we can further bound the endpoints of the $[\Smin,$ $\Smax]$ interval:
\begin{align}\nonumber
\Smin = I^{\amin}_0 &\leq I^{\amin}_0 + \amin(c_1+c_2-r)-\amin^2-c_2(n-1-r) \\ \nonumber
    &=\frac{r^2}{2}+ \amin(c_1-c_2+r-1)+c_2(n-r-1) -r\\\nonumber
    &\leq \amin(c_1-c_2) +c_2 n + r\left(\amin - c_2 + \frac{r}{2} \right) \\ \label{eq:sminlower}  &\leq (\amin+r)(c_1-c_2) + nc_2 \\
\nonumber
\Smax =I^{\amax}_1 &\geq I^{\amax}_1 + \amax^2-\amax(c_1+c_2-r) + c_2(n-1-r) \\ \nonumber
    &=\frac{r^2}{2} + \amax(c_1-c_2+r+1)+c_2(n-1-r)+r \\ \nonumber
    &\geq \amax(c_1 - c_2) + c_2 n + r\left(\amax - c_2 + \frac{r}{2} \right)\\ \label{eq:smaxlower}
    &\geq \amax(c_1 - c_2) + c_2 n
\end{align}
Thus, to show $\Smin \leq \Sigma \leq \Smax$ it suffices to prove
\[ \amin +r \leq \frac{\Sigma - n c_2}{(c_1-c_2) } \leq \amax \]
From \eqref{eq:sigma_abs_bound} we know that
\[ (c_1+c_2+1) - (1-\varepsilon)\sqrt{Q} \leq \frac{\Sigma-nc_2}{(c_1-c_2)/2} \leq (c_1+c_2+1)+(1-\ep)\sqrt{Q},\] hence we will be done if we choose $\ep$ such that
\begin{align*}
2r+2\amin&\leq(c_1+c_2+1)-(1-\ep)\sqrt{Q} \;\mbox{ and}\\
2\amax &\geq (c_1+c_2+1)+(1-\ep)\sqrt{Q}
\end{align*}
We know that \(\amax\) is the largest integer solution of \eqref{eq:newxbound} hence
\[\amax \geq \frac{(c_2+c_1-r)+\sqrt{(c_2+c_1-r)^2-4c_2(n-r-1)-4r}}{2} - 1\] so
\begin{align*}
     2\amax - (c_1+c_2+1) &\geq \sqrt{(c_2+c_1-r)^2-4(c_2(n-r-1)+r)} - (r+3)\\
     &= \sqrt{ (c_1-c_2 -r)^2 - 4c_2(n-c_1-1)-4r} - (r+3)
\end{align*}
Similarly, $\amin$ is the smallest integer solution of \eqref{eq:newxbound}, so
\[\amin \leq \frac{(c_2+c_1-r)-\sqrt{(c_2+c_1-r)^2-4c_2(n-r-1)-4r}}{2} + 1\] so
\begin{align*}
     2r + 2\amin - (c_1+c_2+1) &\leq -\sqrt{(c_2+c_1-r)^2-4(c_2(n-r-1)+r)} + (r+1)\\
     &= -\sqrt{ (c_1-c_2 -r)^2 - 4c_2(n-c_1-1)-4r} + (r+1)
\end{align*}
To simplify notation, it will be convenient to introduce:
\[ Q(r) := (c_1-c_2-r)^2 - 4c_2(n-1-c_1)+4r.\]
Now we simply set
\begin{equation}\label{eq:epsilon}
\varepsilon = 1 -\frac{\sqrt{Q(r)}-r-3}{\sqrt{Q}}.
\end{equation}
This is equivalent to
\[ (1-\ep)\sqrt{Q} = \sqrt{Q(r)} -(r+3).\]
Putting everything together, for $\Sigma$ satisfying \eqref{eq:sigma_abs_bound} we can combine all the above steps and get
\begin{align*}
\frac{\Smin - c_2 n}{(c_1-c_2)/2} - (c_1+c_2+1)&\leq 2r+\amin -(c_1+c_2+1)\\ &\leq r+1 - \sqrt{Q(r)}\\ &\leq
-(1-\varepsilon)\sqrt{Q}  \\ &\leq\frac{\Sigma - n c_2}{(c_1-c_2)/2}-(c_1+c_2+1) \\ &\leq (1-\varepsilon)\cdot \sqrt{Q} \\ &= \sqrt{Q(r)}-(r+3)\\ &\leq 2\amax - (c_1+c_2+1) \leq \frac{\Smax - n c_2}{(c_1-c_2)/2}-(c_1+c_2+1)
\end{align*}
and hence $\Smin \leq \Sigma \leq \Smax$

The only task left is to bound $\varepsilon$. To that end, let us calculate
\begin{align*}\sqrt{Q}-(\sqrt{Q(r)}-r-3) &=r+3+ \frac{Q-Q(r)}{\sqrt{Q}+\sqrt{Q(r)}}\\
&= r+3+\frac{(r+1)(2c_1-2c_2+1-r)-4r}{\sqrt{Q}+\sqrt{Q(r)}}\\
&\leq (r+3) \cdot \left(\frac{\sqrt{Q}+2c_1-2c_2+1-r}{\sqrt{Q}}\right) \\
&\leq (r+3)\cdot \left(\frac{3c_1 - 3c_2 +2 -r}{\sqrt{Q}}\right)\\
&\leq 3(r+3) \cdot \frac{(c_1-c_2)}{\sqrt{Q}},
\end{align*}
where we used $\sqrt{Q} \leq c_1 - c_2 +1$ along with some rather crude estimates. Recall that our assumption was $(1-\beta)(c_1-c_2+1)^2 \geq 4c_2(n-1-c_1)$, equivalently $Q > \beta (c_1-c_2+1)^2$. Using this assumption on $Q$ we see that
\[ \varepsilon = \frac{\sqrt{Q}-\sqrt{Q(r)}+r+3}{\sqrt{Q}} \leq 3(r+3) \frac{c_1-c_2}{Q} \leq 3(r+3) \frac{c_1-c_2}{\beta (c_1-c_2+1)^2} \leq \frac{3(r+3)}{\beta(c_1-c_2)}\]
as claimed.
\end{proof}

\end{document}
